\providecommand{\Res}[2]{\qty[round-mode=figures,round-precision=3]{\fpeval{#1}}{#2}}
\providecommand{\Val}[1]{\num[round-mode=figures,round-precision=4]{\fpeval{#1}}}
\providecommand{\ResSci}[2]{\qty[round-mode=figures,round-precision=3,exponent-mode=scientific]{\fpeval{#1}}{#2}}
\subsubsection*{Question a) Angle de la frange d'ordre 1}
\[ d\sin\theta_1 = \lambda \quad\Rightarrow\quad \theta_1 = \arcsin\left(\frac{\num{\vsi{lam}}}{\num{\vsi{sep}}}\right) = \Res{asind(\vsi{lam}/\vsi{sep})}{\degree} \]
La frange d'ordre 1 est à \textbf{\Res{asind(\vsi{lam}/\vsi{sep})}{\degree}}.

\subsubsection*{Question b) Intensité de la frange d'ordre 1}
\begin{enumerate}[label=\arabic*)]
    \item \textbf{Variables connues et recherchées :}
    \begin{itemize}[label= -]
        \item $I_0 = \qty{\v{I0}}{W/m^2}$, $a = \qty{\v{larg}}{\micro\metre}$, $d = \qty{\v{sep}}{\micro\metre}$, $m = 1$
        \item Variable recherchée : $I_1$
    \end{itemize}
    \item \textbf{Équations pertinentes :}
    \[ I = I_0\cos^2\left(\frac{\phi}{2}\right)\left(\frac{\sin\alpha}{\alpha}\right)^2, \qquad \alpha = \frac{\pi a\sin\theta}{\lambda} \]
    \item \textbf{Isolation de la variable recherchée :} à une frange brillante, $\cos^2(\phi/2) = 1$ et $\sin\theta_m = m\lambda/d$, donc $\alpha = \dfrac{\pi m a}{d}$ et
    \[ I_m = I_0\left(\frac{\sin(\pi m a/d)}{\pi m a/d}\right)^2 \]
    \item \textbf{Résolution numérique :} $\alpha = \pi \cdot \num{\v{larg}}/\num{\v{sep}} = \Res{pi*\v{larg}/\v{sep}}{rad}$
    \[ I_1 = \num{\v{I0}} \cdot \left(\frac{\sin(\Val{pi*\v{larg}/\v{sep}})}{\Val{pi*\v{larg}/\v{sep}}}\right)^2 = \Res{\v{I0}*(sin(pi*\v{larg}/\v{sep})/(pi*\v{larg}/\v{sep}))^2}{W/m^2} \]
    \item \textbf{Réponse :} la frange d'ordre 1 a une intensité de \textbf{\Res{\v{I0}*(sin(pi*\v{larg}/\v{sep})/(pi*\v{larg}/\v{sep}))^2}{W/m^2}}.
\end{enumerate}

\subsubsection*{Question c) Intensité de la frange d'ordre 2}
Avec $m = 2$, $\alpha = 2\pi a/d = \Res{2*pi*\v{larg}/\v{sep}}{rad}$ :
\[ I_2 = \num{\v{I0}} \cdot \left(\frac{\sin(\Val{2*pi*\v{larg}/\v{sep}})}{\Val{2*pi*\v{larg}/\v{sep}}}\right)^2 = \Res{\v{I0}*(sin(2*pi*\v{larg}/\v{sep})/(2*pi*\v{larg}/\v{sep}))^2}{W/m^2} \]
La frange d'ordre 2 a une intensité de \textbf{\Res{\v{I0}*(sin(2*pi*\v{larg}/\v{sep})/(2*pi*\v{larg}/\v{sep}))^2}{W/m^2}}.
